\documentclass{article}

\usepackage[a4paper, margin=1in]{geometry}
\usepackage[french]{isodate}
\usepackage{setspace}

\doublespacing

\title{Lecture 2}
\author{Aden Perry}
\date{\today}

\begin{document}
\maketitle

\section{Les mots}

\begin{figure}[ht]
	\center
	\begin{tabular}{rl}
		Français    & English       \\
		\hline
		\hline
		également   & also          \\
		accueillir  & accomadate    \\
		fourmilière & activity hive \\
		dispositif  & plan          \\
		orthographe & spelling      \\
		fournir     & supply        \\
		démarrage   & beginnning    \\
		soutenir    & support       \\
		quant à     & as for        \\
		croissance  & growth        \\
		réseau      & network       \\
		bâtiment    & property      \\
		carré       & square        \\
		coûteux     & costly        \\
		\hline
		\hline
	\end{tabular}
\end{figure}

\section{Les questions}

\begin{enumerate}
	\item Quelle est la différence principale entre un incubateur et un
	      accélérateur de startups?

	      Les incubaateurs créent des biz ; les accélérateurs les développent.

	\item Qui est le fondateur de Station F et en quelle année ce projet a-t-il
	      vu le jour?

	      Xavier Niel a fondé ce projet. Il at ouvert ses portes en 2017.

	\item Combien de startups Station F accueille-t-il actuellement?

	      Station F accueille plus de 1.000 startups aujourd'hui.

	\item Citez deux écoles de commerce qui ont ouvert un incubateur à Station
	      F.

	      HEC Paris et l'EDHEC ont ouvert un incubateur à Station F.

	\item Nommez trois grands groupes qui ont déployé des programmes ou des
	      incubateurs au sein de Station F.

	      Ubisoft, Microsoft, et Slack ont ont déployé des programmes ou des
	      incubateurs au sein de Station F.

	\item Quel est l'objectif du programme « Fighters Program » proposé par
	      Station F?

	      L'objectif du programme « Fighters Program » proposé par Station F
	      est d'aider aux entrepreneurs issus de milieux défavorisés.

	\item Comment peut-on intégrer Station F si on a un projet de startup?

	      Si on a un projet de startup, on peut postuler aux programmes
	      proposés par les écoles de commerce ou díngenieurs.

	\item Quel grand restaurant réputé s’est installé à Station F et combien de
	      personnes peut-il accueillir?

	      Big Mamma est installé á Station F. Il peut accueillir jusqu'à 1.000
	      personnes.

\end{enumerate}

\section{La recherche}

\begin{enumerate}
	\item \textbf{Mithril Security:}

          Il est un startup securité qui utilise IA. Pour transparence, se
          produit est open source. Théoriquement, l'administration de Mithril
          ne peut pas accéder à l'informatique de ses clients.

	\item \textbf{Hugging Face}


	\item \textbf{BigBlue:}


\end{enumerate}

\end{document}
